\section{The Justice Criterion: Genuine and Forged Trust}
\label{p13:sec:justice}

\noindent
The mechanism of \xref{p13:sec:mechanism} furnished, as a by-product, a mechanical characterisation of its own counterfeit: forged trust as a precision injected from without rather than earned by a real self-committing likelihood, and the geometric reading of \xref{p13:sec:holonomy} gave the same counterfeit its geometric form, as a phase pumped into a driven path rather than accumulated by an adiabatic loop. But the mechanical and the geometric characterisations are, by themselves, value-neutral. They tell us that the forged case is unstable\,---\,that a precision unsupported by likelihood will be overtaken by accumulating error, that a phase pumped from outside is not the loop's own\,---\,but instability is not yet injustice, and a transient is not yet a wrong. This section makes the further claim: that the line between genuine and forged trust, which the mechanism drew as a line between the stable and the unstable, is also and irreducibly a line of justice, between a way of producing trust that honours the one in whom it is produced and a way that wrongs him. The claim is that these two lines\,---\,the mechanical and the ethical\,---\,coincide, and the section's task is to say why they coincide without pretending that either is derivable from the other.

\subsection{Why the ethical criterion is not derivable from the mechanism}
\label{p13:sub:just-notderivable}

It must be said first what the section does \emph{not} claim, for the series' discipline requires it and the temptation to overreach is real. It does not claim that the wrongness of forged trust follows from its instability\,---\,that the forged case is unjust \emph{because} it is destined to collapse. That would be to derive an ethical conclusion from a dynamical premise, to read a wrong off a trajectory, and the series has been at pains throughout to keep these orthogonal. The mechanism answers a dynamical question: will this trust survive, is it well-founded, will the error signals confirm or refute it. The justice criterion answers a different question, which no amount of dynamical information settles: is this a way of treating the other that honours him or wrongs him. A forged trust that happened, against the odds, to be stable\,---\,a manipulation that was never found out, a pumped phase that never collapsed\,---\,would be no less a wrong for its stability; and a genuine trust that collapsed through no one's fault\,---\,a real self-commitment overtaken by circumstance\,---\,would be no less honourable for its collapse. The two criteria can be made to agree, and the section will argue that at the limit they do; but they are not the same criterion, and the agreement is a substantive finding rather than a definitional identity. This is the orthogonality the series has insisted upon elsewhere, between the question of whether a cycle lives and the question of whether it is just, and it holds here in the form: the mechanism tells whether the trust stands, the justice criterion tells whether it should have been produced as it was.

\subsection{The single line, drawn ethically}
\label{p13:sub:just-line}

With that established, the coincidence can be drawn. The mechanical mark of manipulation was stated in \xref{p13:sub:mech-forged}: to induce the other's generative model to lower its prediction error without supplying the high-fidelity likelihood that would warrant the lowering. The ethical mark of manipulation is the same act described under a different aspect, and the section names three aspects of the one act, each of which is a way the manipulation wrongs the one manipulated.

The first aspect is epistemic. To induce in another a high-precision prediction unsupported by any real likelihood is to corrupt his capacity to predict the world truly\,---\,to install in his model a confidence that the world will not bear out, and so to set him at odds with his own future experience. The manipulated person is not merely deceived about a particular fact; his instrument for forming expectations has been tampered with, made to deliver a precision the evidence does not warrant, so that he is rendered, in just this respect, a worse knower than he was. Manipulation is an injury to the other's epistemic autonomy: it does not persuade him by giving him grounds he can weigh, which would leave his judgment intact and free; it works on his model from outside, behind the back of his judgment, to produce in him an expectation he did not form on grounds and could not have refused. This is the precise epistemic sense in which manipulation differs from honest persuasion, and it is a difference the mechanism makes exact: persuasion supplies a likelihood the other's model may weigh and accept or reject; manipulation injects a precision the model did not earn and cannot examine.

The second aspect is the Kantian, and it returns the paper to the distinction of \xref{p13:sec:freedom}. To produce trust by injecting precision rather than by supplying a real self-committing likelihood is to treat the other not as a free mind forming its own expectations on grounds but as a mechanism whose states are to be arranged from outside\,---\,to apply to him the mechanical cause rather than the inferential, to manage him rather than to address him. It is, in the terms of the second formulation of the categorical imperative, to treat him merely as a means: as an object whose predictive states one engineers toward one's own ends, rather than as an end in himself who forms his expectations freely and could withhold them~\parencite{kant1998}. The genuine production of trust, by contrast, addresses the other exactly as \xref{p13:sec:freedom} required: it supplies him with a real ground\,---\,a costly, irrevocable self-commitment\,---\,which his free judgment may take up, and it leaves the taking-up to him. The difference between genuine and forged trust is thus, in its second aspect, the difference between addressing a free being as free and handling a free being as a mechanism; and this is why the forged case is not merely a technical defect in the production of trust but a wrong done to the one in whom it is produced.

The third aspect is the one the series has carried from its work on generative justice, and it is the deepest, for it names what the manipulator does with the relation as a whole. To manipulate is to treat the other's trust as \emph{fuel}\,---\,to draw upon his capacity to rely, his openness to being committed-to, as a resource to be consumed for an end that is not his and from which he is excluded. The forged relation is one in which the trusting party is rendered a mere means to a value that returns to the manipulator and not to him: his trust is mined, his vulnerability is exploited, and he is positioned in the cycle as the one who supplies what others consume. This is the criterion the series has stated as the question to be put to any cycle that deepens: is there someone within it who figures only as a cost, excluded from the return\,---\,is there one who is rendered mere fuel? Forged trust answers yes; it is the production of an expectation whose function is to extract from the one who holds it, and who is, in the holding, used up. Genuine trust answers no; it is the production of an expectation that opens a relation in which the trusting party is a beneficiary and not merely a source, in which what is generated returns to the one whose trust made it possible. The justice of trust, in this third and deepest aspect, is the justice of the return: genuine trust returns to the one who gives it, forged trust extracts from him.

\subsection{The three aspects are one}
\label{p13:sub:just-oneact}

These three\,---\,the epistemic injury, the Kantian reduction, and the rendering-as-fuel\,---\,are not three separate wrongs that manipulation happens to commit together; they are three aspects of a single act, the act the mechanism described as the injection of unwarranted precision. To inject precision without likelihood just \emph{is} to corrupt the other's knowing (the epistemic aspect), which just is to bypass his free judgment and handle him as a mechanism (the Kantian aspect), which just is to position him as a source to be drawn upon rather than a partner in a return (the aspect of fuel). The mechanism gave the act its mechanical description; the three aspects give the one act its ethical description; and that the mechanical and the ethical descriptions are descriptions of the same act is why the line the mechanism drew and the line justice draws coincide. They coincide not because the one is derived from the other\,---\,the section began by refusing that\,---\,but because they are two descriptions, at two levels, of one and the same thing: the producing of an expectation in another without the real self-commitment that would honour him in the producing of it.

\subsection{The criterion stated, and what it leaves to practice}
\label{p13:sub:just-criterion}

The criterion may now be stated in a form the rest of the paper will use. \emb{Genuine trust is an expectation produced in another by supplying a real ground his free judgment may take up, and which opens a relation that returns to him; forged trust is an expectation induced in another by injecting a precision his judgment cannot examine, which handles him as a mechanism and draws upon him as fuel.} The criterion is, as the series requires, a criterion of justice laid alongside the criterion of dynamics, not derived from it: the mechanism tells whether an expectation will stand, the criterion tells whether it was rightly produced, and the two must each be applied, neither collapsing into the other.

One consequence must be marked, for it sets the problem the praxis will have to solve and guards against a misreading the rest of the paper cannot afford. The genuine and the forged are not, as outward performances, reliably distinguishable. The manipulator who stages a costly-looking commitment and the sincere promisor who makes a real one may, in the moment, present indistinguishable surfaces; the injected precision and the earned precision may, for a time, feel the same from inside. What distinguishes them is not a difference in the performance but a difference in what lies behind it\,---\,whether there is, in fact, a real self-commitment and a real return, or only their appearance. This is why the line, though sharp in principle, is treacherous in practice, and why the production of genuine trust cannot be reduced to the production of its signs: the signs can be forged, and the difference that matters is precisely the one the signs do not, by themselves, reveal. The translation of sincerity across the borders of epistemology (\xref{p13:sec:translation}) and the praxis of trust (\xref{p13:sec:praxis}) are, in large part, the working-out of how one acts well under exactly this condition\,---\,that the genuine must be produced and offered in a world where its signs are indistinguishable from the forged, and where the difference, though real, is not on the surface to be read. The criterion of justice tells us what the difference is; it falls to practice to live by a difference one cannot, from the outside, simply see.
